\paragrafos{Η έννοια της εξίσωσης: ο αριθμός που κρύβεται}{Α.4.1, «Η έννοια της εξίσωσης – Οι εξισώσεις $\alpha+x=\beta$, $x-\alpha=\beta$, $\alpha-x=\beta$, $\alpha x=\beta$, $\alpha:x=\beta$, $x:\alpha=\beta$»}

\begin{anagke}
Στο διάλειμμα η Πολυξένη, που αγαπάει τους κωδικούς, προτείνει ένα παιχνίδι: «Σκέψου έναν αριθμό, πολλαπλασίασέ τον επί τρία, πρόσθεσε πέντε και πες μου μόνο το αποτέλεσμα. Εγώ θα βρω τον αριθμό σου». Ο Νίκος σκέφτεται και λέει «είκοσι έξι». Η Πολυξένη απαντάει αμέσως: «Επτά!». Ο Νίκος ανοίγει το στόμα: ήταν πράγματι το επτά.

Η Μυρτώ δοκιμάζει με το $41$ και η Πολυξένη βρίσκει πάλι τον αριθμό. «Μαγεία;» ρωτάει ο Στέλιος. «Όχι», απαντάει η Πολυξένη. «Αναποδογυρίζω τα βήματα: αφαιρώ το πέντε και διαιρώ με το τρία». Ο Βασίλης αμφιβάλλει: «Μα εσύ δεν ξέρεις τον αριθμό. Πώς γράφεις κάτι που δεν ξέρεις;». Η Αγγελική παίρνει το μολύβι και ζωγραφίζει στο τετράδιο ένα κουτάκι: «Βάζω ένα κουτάκι εκεί που κρύβεται ο αριθμός. Τρία κουτάκια συν πέντε κάνουν είκοσι έξι».

Στο μάθημα της κυρίας Αλεξάνδρας η Φανή ρωτάει: «Αν αντί για κουτάκι γράψουμε ένα γράμμα, το ίδιο παιχνίδι λειτουργεί για κάθε αριθμό;». Εκεί αρχίζουν τα γράμματα στα μαθηματικά. Το γράμμα είναι ένα κουτάκι για τον αριθμό που δεν ξέρουμε ακόμα, και με τη βοήθειά του μπορούμε να γράψουμε το παιχνίδι ως ισότητα και να το λύσουμε.
\end{anagke}

\begin{lexi}{εξίσωση, άγνωστος, άλγεβρα}{εξίσωση, άγνωστος, ρίζα}
Η λέξη \lexiword{εξίσωση} βγαίνει από το «ίσος»: εξισώνω σημαίνει «κάνω ίσα» δύο ποσά, όπως οι δύο δίσκοι μιας ζυγαριάς. \lexiword{Άγνωστος} είναι ο αριθμός που δεν ξέρουμε και ψάχνουμε. Η \lexiword{ρίζα} της εξίσωσης είναι η λύση της, και λέγεται έτσι γιατί «βγαίνει από το βάθος», όπως η ρίζα ενός φυτού. Η λέξη \lexiword{άλγεβρα} έρχεται από το αραβικό «αλ-τζαμπρ» (αποκατάσταση), από τον τίτλο ενός βιβλίου του Αλ-Χουαρίζμι τον 9ο αιώνα μ.Χ., όπου λύνονταν εξισώσεις.
\end{lexi}

\begin{metafrasi}
Ο επόμενος και ο προηγούμενος του $\nu$. & $\nu+1$ και $\nu-1$\\
Ένας άρτιος και ένας περιττός. & $2\nu$ και $2\nu+1$\\
Τα πολλαπλάσια του $3$. & $3\nu$\\
Το διπλάσιο του $x$. & $2x$\\
Ένας αριθμός αυξημένος κατά $8$, ελαττωμένος κατά $4$. & $x+8$ και $x-4$\\
Το τετραπλάσιο ενός αριθμού, αυξημένο κατά $2$, κάνει $22$. & $4x+2=22$\\
Αν σε έναν αριθμό προσθέσω $5$ βγαίνει $8$. & $x+5=8$\\
\end{metafrasi}

\tmima{Η θεωρία}

\tmimasmall{Γράμματα στη θέση αριθμών}

Με ένα γράμμα (π.χ. $x$, $\alpha$, $\nu$) παριστάνουμε έναν αριθμό, γνωστό ή άγνωστο, και με γράμματα και πράξεις γράφουμε προτάσεις συνοπτικά. Αν η πλευρά ενός τετραγώνου είναι $\alpha$, η περίμετρός του είναι $4\alpha$ και το εμβαδόν του $\alpha\cdot\alpha$. Μια έκφραση με γράμματα και αριθμούς λέγεται \textbf{αλγεβρική παράσταση}.

Όταν ο ίδιος άγνωστος εμφανίζεται πολλές φορές, συντομεύουμε με την επιμεριστική ιδιότητα: $x+x+x+x=4x$, $3\alpha+5\alpha=(3+5)\alpha=8\alpha$, $15\beta-9\beta=6\beta$, $18x+7x+4x=29x$. Μπορούμε να προσθέτουμε ή να αφαιρούμε μόνο τους όρους που έχουν το \emph{ίδιο γράμμα}: $4\alpha+2\beta+\alpha+3\beta=5\alpha+5\beta$. Δεν συνενώνουμε το $5\alpha$ με το $5\beta$.

\tmimasmall{Τι είναι εξίσωση}

Η ισότητα $2\cdot6=12$ μπορούμε μόνο να τη βεβαιώσουμε. Η ισότητα $2\cdot x=12$ δεν είναι ίδια: κρύβει έναν αριθμό που, αν τον βάλουμε στη θέση του $x$, την κάνει αληθινή, ενώ κάθε άλλος αριθμός την κάνει ψευδή.

\begin{orismos}{εξίσωση με έναν άγνωστο}
\textbf{Εξίσωση} με έναν άγνωστο λέγεται μια ισότητα που περιέχει αριθμούς και ένα γράμμα (τον άγνωστο), π.χ. $x+5=12$, $10-z=1$, $7\varphi=14$, $24:\psi=6$.\\
\textbf{Λύση ή ρίζα} της εξίσωσης είναι ο αριθμός που, όταν αντικαταστήσει τον άγνωστο, \textbf{επαληθεύει} την ισότητα. Η διαδικασία εύρεσης της λύσης λέγεται \textbf{επίλυση}.
\end{orismos}

Για παράδειγμα η ρίζα της $x-7=5$ είναι το $12$, γιατί $12-7=5$. Γράφουμε $x=12$.

\tmimasmall{Ταυτότητα και αδύνατη εξίσωση}

\begin{itemize}[leftmargin=1.2em,itemsep=0.3ex,topsep=2pt]
\item Μια εξίσωση λέγεται \textbf{ταυτότητα} (ή αόριστη) όταν \emph{κάθε} αριθμός είναι λύση της, π.χ. $x=x$ ή $0\cdot z=0$.
\item Λέγεται \textbf{αδύνατη} όταν \emph{κανένας} αριθμός δεν την επαληθεύει, π.χ. $x+2=x+6$ ή $0\cdot\omega=5$.
\end{itemize}

\tmimasmall{Οι έξι βασικές εξισώσεις}

Από τον ορισμό των πράξεων (η αφαίρεση είναι η αντίστροφη πράξη της πρόσθεσης και η διαίρεση του πολλαπλασιασμού) βρίσκουμε τη λύση κάθε μίας:

\begin{center}
\renewcommand{\arraystretch}{1.5}
\begin{tabular}{@{}lll@{}}
\toprule
\textbf{Εξίσωση} & \textbf{Λύση} & \textbf{Παράδειγμα}\\
\midrule
$x+\alpha=\beta$ & $x=\beta-\alpha$ & $x+5=12\ \Rightarrow\ x=12-5=7$\\
$x-\alpha=\beta$ & $x=\beta+\alpha$ & $y-2=3\ \Rightarrow\ y=3+2=5$\\
$\alpha-x=\beta$ & $x=\alpha-\beta$ & $10-z=1\ \Rightarrow\ z=10-1=9$\\
$\alpha\cdot x=\beta$ & $x=\beta:\alpha$ & $7\varphi=14\ \Rightarrow\ \varphi=14:7=2$\\
$x:\alpha=\beta$ & $x=\beta\cdot\alpha$ & $\omega:5=4\ \Rightarrow\ \omega=4\cdot5=20$\\
$\alpha:x=\beta$ & $x=\alpha:\beta$ & $24:\psi=6\ \Rightarrow\ \psi=24:6=4$\\
\bottomrule
\end{tabular}
\end{center}

(Στις δύο τελευταίες και στην τέταρτη προϋποθέτουμε ότι δεν διαιρούμε με το $0$.) Μετά τη λύση \textbf{επαληθεύουμε}: βάζουμε την τιμή στην αρχική εξίσωση και ελέγχουμε ότι βγαίνει ισότητα. Στις εξισώσεις δεν γράφουμε τις μονάδες, αλλά πρέπει να ξέρουμε σε ποιες μονάδες δουλεύουμε.

\begin{prosoxi}
\begin{itemize}[leftmargin=1.2em,itemsep=0.3ex,topsep=2pt]
\item Στην $\alpha-x=\beta$ η λύση είναι $\alpha-\beta$, όχι $\beta-\alpha$. Η $12-x=5$ δίνει $x=7$, όχι $17$.
\item Στις $\alpha:x=\beta$ και $x:\alpha=\beta$ οι λύσεις διαφέρουν: $24:\psi=6$ δίνει $\psi=24:6=4$, ενώ $\omega:6=24$ δίνει $\omega=24\cdot6=144$.
\item Το $3\alpha+5\beta$ δεν γίνεται $8\alpha\beta$ ούτε $8\alpha$: έχουμε διαφορετικά γράμματα και δεν συντομεύεται.
\item Δεν αντικαθιστούμε τον άγνωστο και μετά «επαληθεύουμε» με την ίδια πράξη που κάναμε: επαληθεύω βάζοντας τη λύση στην \emph{αρχική} εξίσωση.
\end{itemize}
\end{prosoxi}

\tmima{Λυμένα παραδείγματα}

\paradeigma{Μεταφράζω σε γλώσσα γραμμάτων}
\tixero Ο Νίκος έχει $x$ ευρώ. Η Αγγελική έχει $7$ ευρώ περισσότερα και μαζί έχουν $25$ ευρώ.\par
\tithelo Να γράψω με γράμματα τα χρήματα της Αγγελικής και την ισότητα για το σύνολο.\par
\possindeo Η Αγγελική έχει $x+7$. Μαζί έχουν $x+(x+7)=2x+7$.\par
\apantisi Αγγελική: $x+7$. Σύνολο: $2x+7=25$.\par
\elegxos Αν ο Νίκος είχε $9$ €, η Αγγελική θα είχε $16$ € και μαζί $25$. Άρα $x=9$ ταιριάζει (με την $2x+7=25$ θα το δούμε στην επόμενη παράγραφο).

\paradeigma{Συντομεύω παραστάσεις}
\tixero Οι παραστάσεις $5x+3x-2x$ και $4\alpha+2\beta+\alpha+3\beta$.\par
\tithelo Να τις γράψω απλούστερα.\par
\possindeo Προσθέτω ή αφαιρώ τους όρους που έχουν το ίδιο γράμμα. $5x+3x-2x=(5+3-2)x=6x$. Και $4\alpha+\alpha=5\alpha$, $2\beta+3\beta=5\beta$.\par
\apantisi $6x$ και $5\alpha+5\beta$.\par
\elegxos Με $x=2$: $10+6-4=12=6\cdot2$. Με $\alpha=1$, $\beta=2$: $4+4+1+6=15=5+10$. Σωστά.

\paradeigma{Ποιος αριθμός επαληθεύει;}
\tixero Η εξίσωση $x-4=1$ και οι αριθμοί $3$, $5$, $8$.\par
\tithelo Ποιος είναι ρίζα.\par
\possindeo Αντικαθιστώ: για $x=3$: $3-4$, δεν γίνεται στους φυσικούς. Για $x=5$: $5-4=1$ σωστό. Για $x=8$: $8-4=4\neq1$.\par
\apantisi Ρίζα είναι ο $5$.\par
\elegxos Από τον κανόνα: $x=1+4=5$. Συμφωνεί.

\paradeigma{Λύνω τις έξι βασικές}
\tixero Οι εξισώσεις $x+7{,}5=12$, $x-3{,}2=5$, $15-x=6{,}4$, $4x=18$, $x:6=2{,}5$, $36:x=9$.\par
\tithelo Να τις λύσω και να επαληθεύσω.\par
\possindeo $x=12-7{,}5=4{,}5$. $x=5+3{,}2=8{,}2$. $x=15-6{,}4=8{,}6$. $x=18:4=4{,}5$. $x=2{,}5\cdot6=15$. $x=36:9=4$.\par
\apantisi $4{,}5$, $8{,}2$, $8{,}6$, $4{,}5$, $15$, $4$.\par
\elegxos $4{,}5+7{,}5=12$, $8{,}2-3{,}2=5$, $15-8{,}6=6{,}4$, $4\cdot4{,}5=18$, $15:6=2{,}5$, $36:4=9$. Όλα σωστά.

\paradeigma{Ταυτότητα ή αδύνατη;}
\tixero Οι εξισώσεις $x+3=3+x$, $x+2=x+6$ και $2x-2x=5$.\par
\tithelo Τι είδους είναι.\par
\possindeo Η πρώτη ισχύει για κάθε αριθμό (αντιμεταθετική ιδιότητα): ταυτότητα. Στη δεύτερη, όποιον αριθμό και να βάλω, το αριστερό μέλος είναι κατά $4$ μικρότερο από το δεξί: αδύνατη. Στην τρίτη το $2x-2x=0$ και δεν ισούται με $5$: αδύνατη.\par
\apantisi Ταυτότητα, αδύνατη, αδύνατη.\par
\elegxos Με $x=10$: $13=13$, αλλά $12\neq16$ και $0\neq5$.

\paradeigma{Το ύψος της δεξαμενής}
\tixero Η νέα δεξαμενή του χωριού έχει όγκο $7{,}2$ m$^3$, μήκος $2$ μ. και πλάτος $1{,}2$ μ.\par
\tithelo Το ύψος της.\par
\possindeo Αν $x$ το ύψος σε μέτρα, ο όγκος είναι $2\cdot1{,}2\cdot x=2{,}4x$. Άρα $2{,}4x=7{,}2$ και $x=7{,}2:2{,}4=3$.\par
\apantisi $3$ μ.\par
\elegxos $2\cdot1{,}2\cdot3=7{,}2$ m$^3$. Σωστό.

\tmima{Ασκήσεις}

\zoni{Το κατάλαβα}{γράμματα, παραστάσεις, έξι βασικές εξισώσεις}

\ex Γράψε με γράμματα (με $x$ έναν αριθμό): α) ο επόμενός του \quad β) ο προηγούμενός του \quad γ) το τριπλάσιό του \quad δ) αυξημένος κατά $12$ \quad ε) ελαττωμένος κατά $5$ \quad στ) το μισό του \quad ζ) το άθροισμά του με το διπλάσιό του.

\ex Γράψε με γράμματα: α) ένας άρτιος φυσικός \quad β) ένας περιττός φυσικός \quad γ) ένα πολλαπλάσιο του $5$ \quad δ) τρεις διαδοχικοί φυσικοί.

\ex Διατύπωσε με λόγια: α) $5x-3$ \quad β) $2(x+1)=14$ \quad γ) $\alpha+\beta$ \quad δ) $6\kappa-\kappa=30$.

\ex Η πλευρά ενός τετραγώνου είναι $\alpha$. Γράψε την περίμετρο και το εμβαδόν του. Κάνε το ίδιο για ένα ορθογώνιο με πλευρές $\alpha$ και $\beta$.

\ex Γράψε απλούστερα: α) $x+x+x$ \quad β) $3\alpha+5\alpha$ \quad γ) $7\beta-2\beta$ \quad δ) $2x+5y+4x-y$ \quad ε) $9\omega-4\omega+\omega$ \quad στ) $10\kappa+3\lambda-10\kappa$.

\ex Εξέτασε αν ο $12$ είναι λύση της $x-5=7$. Ποιος από τους $2,\,4,\,6$ είναι λύση της $3x=12$;

\ex Βρες τη λύση: α) $x+6=10$ \quad β) $x-2=7$ \quad γ) $15-x=9$ \quad δ) $4x=20$ \quad ε) $x:3=5$ \quad στ) $18:x=6$.

\ex Λύσε: α) $x+2{,}75=6$ \quad β) $x-0{,}8=1{,}15$ \quad γ) $12-x=4{,}35$ \quad δ) $0{,}5x=3{,}2$ \quad ε) $x:0{,}4=7$ \quad στ) $9:x=0{,}6$.

\ex Λύσε: α) $x+\frac{1}{4}=\frac{3}{4}$ \quad β) $\frac{2}{3}x=\frac{4}{9}$ \quad γ) $\frac{5}{6}-x=\frac{1}{3}$ \quad δ) $x:3=\frac{2}{9}$.

\ex Σωστό ή λάθος; α) Η $3+x=10$ έχει λύση το $7$. \quad β) Η $x-x=0$ είναι ταυτότητα. \quad γ) Η $x+1=x+2$ είναι αδύνατη. \quad δ) Η $0\cdot x=5$ είναι αδύνατη. \quad ε) Η $5x=0$ έχει λύση το $0$. \quad στ) Η $6:x=2$ έχει λύση το $12$.

\zoni{Το χρησιμοποιώ}{ζυγαριές, δεξαμενές, μυστικοί αριθμοί}

\ex \textbf{Η ζυγαριά.} Μια ζυγαριά ισορροπεί όταν στον έναν δίσκο υπάρχει μια σοκολάτα βάρους $x$ g και στον άλλον $150$ g και μισή σοκολάτα. Γράψε την ισότητα και βρες το βάρος της σοκολάτας με σκέψη.

\ex \textbf{Η δεξαμενή.} Μια δεξαμενή έχει όγκο $12$ m$^3$, μήκος $2$ μ. και πλάτος $1{,}5$ μ. Βρες το ύψος της με εξίσωση.

\ex \textbf{Ο μυστικός αριθμός.} Η Πολυξένη λέει: «Σκέφτηκα έναν αριθμό και βρήκα...». Γράψε εξίσωση και βρες τον αριθμό όταν: α) του πρόσθεσε $17$ και βρήκε $52$ \quad β) αφαίρεσε $9$ και βρήκε $14$ \quad γ) τον αφαίρεσε από το $50$ και βρήκε $18$ \quad δ) τον πολλαπλασίασε επί $6$ και βρήκε $84$ \quad ε) τον διαίρεσε με $8$ και βρήκε $7$.

\ex \textbf{Τα εισιτήρια.} Τέσσερα ίδια εισιτήρια πλοίου κοστίζουν $54$ €. Πόσο κάνει το καθένα; Γράψε την εξίσωση.

\ex \textbf{Ο κουμπαράς.} Ο Ηλίας έχει $36$ € και ξοδεύει κάποιο ποσό. Του μένουν $11{,}5$ €. Πόσα ξόδεψε;

\ex \textbf{Το μοτίβο.} Βάζουμε σε σειρά τετράγωνα πλευράς $1$ cm: το πρώτο σχήμα έχει $1$ τετράγωνο (περίμετρος $4$ cm), το δεύτερο $2$ (περίμετρος $6$ cm), το τρίτο $3$ (περίμετρος $8$ cm). α) Πόση είναι η περίμετρος του πέμπτου σχήματος; β) Γράψε έναν τύπο για το $\nu$-οστό σχήμα. γ) Ποιο σχήμα έχει περίμετρο $98$ cm;

\zoni{Το σκέφτομαι}{όταν το γράμμα κρύβει κάτι}

\ex \textbf{Το κόλπο της Πολυξένης.} Σκέψου έναν αριθμό $x$, πρόσθεσε $4$, διπλασίασε και αφαίρεσε $8$. Γράψε με παράσταση τι κάνεις και δείξε ότι το αποτέλεσμα είναι πάντα το διπλάσιο του αριθμού που σκέφτηκες.

\ex Η εξίσωση $x+8=5$ δεν έχει λύση στους φυσικούς αριθμούς. Γιατί; Τι είδους αριθμός θα την έλυνε;

\ex Ο Στέλιος λύνει την $12-x=5$ και γράφει $x=12+5=17$. Βρες το λάθος, το σωστό και επαλήθευσε.

\ex Γιατί η $0\cdot x=0$ έχει άπειρες λύσεις και η $0\cdot x=5$ καμία; Εξήγησε με το «ποιος αριθμός επί $0$ δίνει...».

\ex Γράψε μια εξίσωση με λύση το $6$, μια αδύνατη και μια ταυτότητα.

\zoni{Για το σπίτι}{φτιάξε το δικό σου}

\ex Ο παππούς λέει: «Σκέφτηκα έναν αριθμό και το τριπλάσιό του είναι $2{,}4$». Ποιος είναι; Γράψε την εξίσωση.

\ex Φτιάξε μια καθημερινή ιστορία που να οδηγεί στην εξίσωση $x+12=30$ και ζήτησε από ένα μέλος της οικογένειας να τη λύσει.
